% from ega4-18-fr.tex
\oldpage[IV-4]{115}

que $B$ est une $A$-algèbre non ramifiée signifie alors (17.4.2) que $\Spec((B\otimes_A B)/\mathfrak J_{B/A})$ est un sous-schéma induit sur une partie \emph{ouverte et fermée} de $\Spec(B\otimes_A B)$ et on sait que pour qu'il en soit ainsi, il faut et il suffit que $\mathfrak J_{B/A}$ soit un idéal \emph{facteur direct} de $B\otimes_A B$ (Bourbaki, \emph{Alg. comm.}, chap. II, \textsection~4, n\textsuperscript{o}~3, prop. 15); mais il revient au même de dire que le $(B\otimes_A B)$-module quotient $(B\otimes_A B)/\mathfrak J_{B/A}$ est projectif (Bourbaki, \emph{Alg.}, chap. II, 3\textsuperscript{e} éd., \textsection~2, n\textsuperscript{o}~2, prop. 4).

(ii) Si on se rappelle qu'un $A$-module plat de présentation finie est projectif et réciproquement (Bourbaki, \emph{Alg. comm.}, chap. II, \textsection~5, n\textsuperscript{o}~2, cor. 2 du th. 1), l'assertion de (ii) résulte de celle de (i) et de (17.6.2).

\begin{proposition}[18.3.2]
\label{IV.18.3.2}
Soient $A$ un anneau, $\mathfrak J$ un idéal de $A$ tel que, pour la topologie $\mathfrak J$-pré-adique, $A$ soit séparé et complet; on pose $A_0=A/\mathfrak J$. Alors le foncteur
\[
  B\longmapsto B\otimes_A A_0
\]
est une équivalence de la catégorie des $A$-algèbres finies et étales, dans la catégorie des $A_0$-algèbres finies et étales.
\end{proposition}

Nous démontrerons d'abord le lemme suivant:

\begin{lemma}[18.3.2.1]
\label{IV.18.3.2.1}
Soient $A$ un anneau, $\mathfrak J$ un idéal de $A$ tel que, pour la topologie $\mathfrak J$-pré-adique, $A$ soit séparé et complet.
\begin{enumerate}
  \item[(i)] Tout $A$-module projectif de type fini $M$ est séparé et complet pour la topologie $\mathfrak J$-pré-adique, donc limite projective des $(A/\mathfrak J^{n+1})$-modules projectifs $M/\mathfrak J^{n+1}M$.
  \item[(ii)] Inversement, posons $A_n=A/\mathfrak J^{n+1}$, et soit $(M_n)$ un système projectif de $A_n$-modules, tel que, pour tout $n$, l'homomorphisme $M_{n+1}\otimes_{A_{n+1}}A_n\to M_n$ déduit du di-homomorphisme de transition $M_{n+1}\to M_n$ soit bijectif. Supposons en outre les $M_n$ projectifs et $M_0$ de type fini. Alors $M=\varprojlim M_n$ est un $A$-module projectif de type fini tel que l'homomorphisme canonique $M\otimes_A A_0\to M_0$ soit bijectif.
\end{enumerate}
\end{lemma}

(i) Il existe un $A$-module libre de type fini $L$ tel que $M$ soit isomorphe à un facteur direct de $L$; comme $L$ est séparé pour la topologie $\mathfrak J$-pré-adique, il en est de même de tout sous-module $N$ de $L$ puisque $\mathfrak J^{n+1}N\subset\mathfrak J^{n+1}L$; en particulier $M$ est séparé pour cette topologie, et comme l'homomorphisme surjectif $f:L\to M$ est continu pour la topologie $\mathfrak J$-pré-adique, son noyau $N$ est fermé pour la topologie induite par celle de $L$; puisque $L$ est complet et $f$ un morphisme strict, on en conclut que $M$ est complet (Bourbaki, \emph{Top. gén.}, chap. IX, 2\textsuperscript{e} éd., \textsection~3, n\textsuperscript{o}~1, prop. 4).

(ii) Il résulte du lemme de Nakayama que si $M_0$ est engendré par une famille finie $(x_{i0})$ de $r$ éléments et si pour tout $n$, $x_{in}$ est un élément de $M_n$ dont $x_{i,n-1}$ est l'image dans $M_{n-1}$, alors $(x_{in})$ $(1\leq i\leq r)$ est un système de générateurs de $M_n$ (Bourbaki, \emph{Alg. comm.}, chap. II, \textsection~3, cor. 2 de la prop. 4). Cela étant, pour tout $n$, posons $L_n=A_n^r$; si $(e_{in})_{1\leq i\leq r}$ est la base canonique de $L_n$, soit $u_n:L_n\to M_n$ l'application $A_n$-linéaire telle que $u_n(e_{in})=x_{in}$ pour tout $i$. Par hypothèse on a une suite exacte scindée
\[
  0\longrightarrow N_n\xrightarrow{v_n}L_n\xrightarrow{u_n}M_n\longrightarrow 0
\]

